\scanpage{31}
On en déduit donc un morphisme can. de multiplicités schématiques
\[
(37)\qquad
\diagram{\begin{tikzcd}[column sep=large,row sep=small]
\widehat{\mathcal M}_{[G]}\arrow[r,hook]&\widehat{\mathcal M}_{g,\nu}\\
(\widehat{\mathcal M}_G,\Gamma)\arrow[u,phantom,"{\wr\ \mathrm{déf}}" description] & (\widehat{\mathcal M}_{g,I},\unsure{\mathfrak S_I|\mathfrak G_I})\arrow[u,phantom,"{\wr}" description]
\end{tikzcd}}
,
\]
qui est encore une \uline{immersion}, fermée cette fois-ci, de multiplicités schématiques propres et lisses sur $\mathbb Z$. Ainsi se trouve explicité l'adhérence de $\mathcal M_{[G]}$ dans $\widehat{\mathcal M}_{g,\nu}$, comme sous-multiplicité propre et lisse sur $\operatorname{Spec}\mathbb Z$ -- qui \unsure{paraît} \inserted{a priori} curieux ! En $=$ temps, cette construction permet de décrire de façon pratiquement triviale les types $[G']$ qui sont spécialisation de $[G]$, puisque ce sont ceux qui proviennent de $\widehat{\mathcal M}_G$. Nous allons \deleted{\unsure{y}} revenir là dessus maintenant.
\mspage{30}
